\section*{Chapter \ref{ch:b2:search-and-unemployment} --- Search and Unemployment}

\begin{solution}{exo:b2:search-and-unemployment:1}
% figdata: bev_summary
2013-Q1: unemployment 12.3 per cent, vacancies 1.3, tightness $1.3/12.3 = 0.106$. 2019-Q4: 7.4 and 2.3, tightness 0.311. 2022-Q2: 6.7 and 3.3, tightness 0.493.
\end{solution}

\begin{solution}{exo:b2:search-and-unemployment:2}
% code: exercises/solutions.py steady_u
$0.03/(0.03 + 0.5) = 0.0566$, or 5.66 per cent.
\end{solution}

\begin{solution}{exo:b2:search-and-unemployment:3}
% code: exercises/solutions.py matches
$m = 0.7 \times \sqrt{0.0018} = 0.0297$; $\theta = 0.5$; $f = 0.0297/0.06 = 0.495$; $q = 0.0297/0.03 = 0.990$.
\end{solution}

\begin{solution}{exo:b2:search-and-unemployment:4}
Opening a vacancy costs nothing but its running cost, so as long as one more vacancy has a positive expected value some
firm will open it. Each extra vacancy raises tightness and lowers the filling rate of all vacancies, which lowers the
expected value; entry stops only when it reaches zero.
\end{solution}

\begin{solution}{exo:b2:search-and-unemployment:5}
% code: exercises/solutions.py wage
$w = 0.5 \times 0.4 + 0.5 \times (1 + 1 \times 0.529) = 0.2 + 0.7645 = 0.9645$ (the listing's 0.9647 uses the unrounded tightness). The term $\beta c\theta = 0.2645$ rewards the worker for the hiring cost the firm
would pay to replace him or her: $c\theta$ is the expected hiring cost per unemployed worker.
\end{solution}

\begin{solution}{exo:b2:search-and-unemployment:6}
% code: listings/dmp.py
Tightness falls (0.529 to 0.436), the wage rises (0.9647 to 0.9679), unemployment rises (5.56 to 6.10 per cent).
\end{solution}

\begin{solution}{exo:b2:search-and-unemployment:7}
% ledger: b2:lalive2008
By $209 - 30 = 179$ weeks. Unemployment lengthened by at least $0.09 \times 179 = 16$ weeks for men and $0.32 \times 179 = 57$ weeks for women.
\end{solution}

\begin{solution}{exo:b2:search-and-unemployment:8}
% code: exercises/solutions.py outcome
Tightness 0.224, unemployment 8.30 per cent, wage 0.977. Workers take a larger share of each match's surplus, so the firm's share
falls, firms open fewer vacancies, and unemployment rises; the wage is higher even though the market is slacker.
\end{solution}

\begin{solution}{exo:b2:search-and-unemployment:9}
% figdata: bev_summary
In 2019-Q4 unemployment was 7.4 per cent with vacancies 2.3; in 2025-Q4, 6.3 with 2.2: less unemployment with no more vacancies.
On a fixed downward-sloping curve, lower unemployment would require more vacancies, so the curve has shifted inward. In the model this
comes from a higher matching efficiency $A$ or a lower separation rate $s$.
\end{solution}

\begin{solution}{exo:b2:search-and-unemployment:10}
With $q = A/x$ and $f = Ax$: $(1 - \beta)(p - z)A = cx(r + s + \beta Ax)$, that is $c\beta Ax^2 + c(r + s)x - (1 - \beta)(p - z)A = 0$, with positive root
$x = \dfrac{-c(r + s) + \sqrt{c^2(r + s)^2 + 4c\beta A^2(1 - \beta)(p - z)}}{2c\beta A}$.
\end{solution}

\begin{solution}{exo:b2:search-and-unemployment:11}
From the firm's equations with $V = 0$: $J = (p - w)/(r + s)$ and $c = qJ$. From the worker's, $(r + s)(W - U) = w - rU$ and $rU = z + f(W - U)$. The sharing rule gives $W - U = \tfrac{\beta}{1 - \beta}J$. Then
$w = rU + (r + s)\tfrac{\beta}{1 - \beta}J = rU + \tfrac{\beta}{1 - \beta}(p - w)$, so $w = (1 - \beta)rU + \beta p$. And $rU = z + f\tfrac{\beta}{1 - \beta}J = z + \tfrac{\beta}{1 - \beta}\theta qJ = z + \tfrac{\beta}{1 - \beta}c\theta$. Substituting: $w = (1 - \beta)z + \beta c\theta + \beta p$.
\end{solution}

\begin{solution}{exo:b2:search-and-unemployment:12}
% code: exercises/solutions.py elasticity
Take logarithms of $(1 - \beta)(p - z) = (c/q)(r + s + \beta f)$, with $d\ln q = -\eta\,d\ln\theta$ and $d\ln f = (1 - \eta)d\ln\theta$: $\frac{p}{p - z}d\ln p = \eta\,d\ln\theta + \frac{\beta f(1 - \eta)}{r + s + \beta f}d\ln\theta$, which gives the formula. With $z = 0.4$, $f = 0.509$: $1.667/(0.5 + 0.5 \times 0.2547/0.2887) = 1.77$.
With $z = 0.9$, $f = 0.190$: $10/(0.5 + 0.5 \times 0.0950/0.1290) = 11.5$. The listing's finite 1 per cent changes give 1.8 and 11.6.
\end{solution}

\begin{solution}{pb:b2:search-and-unemployment:1}
% code: exercises/solutions.py BRENLAND, outcome, elasticity
\textbf{Part I.}
1. $\dot u = 0.05(1 - u) - f(\theta)u$.
2. $0.05/(0.05 + 0.45) = 10$ per cent.
3. $1/0.45 = 2.2$ quarters.
4. $1/0.05 = 20$ quarters.
5. More jobs end, so for every vacancy rate more workers are unemployed: at each $\theta$, $u = s/(s + f(\theta))$ is higher.
\textbf{Part II.}
6. $f = 0.6\theta^{0.5}$, $q = 0.6\theta^{-0.5}$.
7. $0.6(\lambda u)^{0.5}(\lambda v)^{0.5} = \lambda \times 0.6u^{0.5}v^{0.5}$.
8. $\sqrt\theta = 0.8$, $\theta = 0.64$.
9. $0.6/0.8 = 0.75$.
10. From $u = 0.05/(0.05 + 0.6\sqrt{v/u})$: $v = (0.05/0.6)^2(1 - u)^2/u = 0.00694(1 - u)^2/u$.
\textbf{Part III.}
11. $w = 1 - 0.06 \times 0.5/q = 1 - 0.05\sqrt\theta$.
12. $w = 0.5 \times 0.6 + 0.5(1 + 0.5\theta) = 0.8 + 0.25\theta$.
13. $0.8 + 0.25x^2 = 1 - 0.05x$, so $x^2 + 0.2x - 0.8 = 0$ and $x = (-0.2 + \sqrt{3.24})/2 = 0.8$.
14. $\theta = 0.64$, $f = 0.48$, $q = 0.75$.
15. $u = 0.05/0.53 = 9.43$ per cent; $v = 0.64 \times 9.43 = 6.04$ per cent.
16. $w = 0.8 + 0.16 = 0.96$ and $1 - 0.05 \times 0.8 = 0.96$.
17. $\theta = 0.676$, $u = 9.21$ per cent.
18. $(1/0.4)/(0.5 + 0.5 \times 0.24/0.30) = 2.5/0.9 = 2.78$: a 1 per cent rise in productivity raises tightness by about 2.8 per cent.
\textbf{Part IV.}
19. The wage curve shifts up, by $0.5 \times 0.1 = 0.05$, to $w = 0.85 + 0.25\theta$.
20. $x^2 + 0.2x - 0.6 = 0$, $x = 0.681$, $\theta = 0.464$, $f = 0.409$.
21. $0.05/(0.05 + 0.409) = 10.9$ per cent.
22. $0.966$, up from 0.960. The higher value of unemployment raises the wage by 0.05 at given tightness, but tightness falls, which lowers the
part of the wage that rewards the firm's hiring cost.
23. Each match leaves the firm a smaller share of the surplus, so fewer vacancies pay for their cost.
24. The search effort of the unemployed: more generous benefits let them search less intensely or wait longer for a
better offer.
25. Brenland's equilibrium has tightness 0.64, unemployment 9.43 per cent and a wage of 0.960; raising the value of unemployment from 0.6 to 0.7 lowers tightness to 0.46,
raises unemployment to 10.9 per cent and the wage to 0.966.
\end{solution}
